% Part: lambda-calculus
% Chapter: syntax
% Section: eta

\documentclass[../../../include/open-logic-section]{subfiles}

\begin{document}

\olfileid[hi]{lam}{syn}{eta}
\olsection{$\eta$-परिवर्तन}

$\lambda$-पदों पर एक और संबंध है। \olref[fv]{sec} में हमने
$\lambd[x][(fx)]$ का उदाहरण उपयोग किया था, जो कोई तर्क स्वीकार करके उस पर
$f$ लागू करता है। दूसरे शब्दों में, यह~$f$ वाला ही फलन है:
$\lambd[x][(fx)]N$ और $fN$ दोनों अपचयित होकर~$fN$ बनते हैं। इस विचार को
निरूपित करने के लिए हम $\eta$-अपचयन (और $\eta$-विस्तार) उपयोग करते हैं।

\begin{defn}[$\eta$-संकुचन, $\eredone$]
  \ollabel{defn:beredone}
  \emph{$\eta$-संकुचन} ($\eredone$) पदों पर वह सबसे छोटा संगत संबंध है जो
  निम्न शर्त संतुष्ट करता है:
  \[
    \lambd[x][M x] \eredone M \text{ बशर्ते } x \notin FV(M)
  \]
\end{defn}

\begin{defn}[$\beta\eta$-अपचयन, $\bered$] \ollabel{defn:bered}
  \emph{$\beta\eta$-अपचयन} ($\bered$) पदों पर वह सबसे छोटा परावर्ती और
  संक्रामक संबंध है जिसमें $\bredone$ और $\eredone$ समाहित हैं; अर्थात्
  परावर्तिता और संक्रामकता के नियमों के साथ निम्न दो नियम:
  \begin{enumerate}
  \item यदि $M \bredone N$, तो $M \bered N$। \ollabel{defn:bered3}
  \item यदि $M \eredone N$, तो $M \bered N$। \ollabel{defn:bered4}
  \end{enumerate}
    
\end{defn}

\begin{defn}
  हम तुल्यता-संबंध $\equal$ को $\eta$-परिवर्तन नियम से विस्तारित करते हैं:
  \[
  \lambd[x][f x] \equal f
  \]
  और विस्तारित संबंध को $\equal[\eta]$ से दर्शाते हैं।
\end{defn}

$\eta$-तुल्यता महत्त्वपूर्ण है, क्योंकि वह लैम्ब्डा पदों की विस्तारिकता से
संबंधित है:

\begin{defn}[विस्तारिकता]
  हम तुल्यता-संबंध $\equal$ को (\ext) नियम से विस्तारित करते हैं:
  \begin{center}
  यदि $Mx \equal Nx$, तो $M \equal N$, बशर्ते $x \notin FV(MN)$।
  \end{center}
  और विस्तारित संबंध को $\equal[\ext]$ से दर्शाते हैं।
\end{defn}

मोटे तौर पर यह नियम कहता है कि फलनों के रूप में देखे गए दो पद समान माने
जाने चाहिए यदि समान तर्क पर उनका व्यवहार समान हो।

अब हम सिद्ध करते हैं कि $\eta$ नियम ठीक विस्तारिकता देता है, और कुछ नहीं।

\begin{thm}
  $M \equal[\ext] N$ तभी और केवल तभी जब $M \equal[\eta] N$।
\end{thm}

\begin{proof}
  पहले हम सिद्ध करते हैं कि $\equal[\eta]$ विस्तारिकता नियम के अधीन संवृत
  है। अर्थात् $ext$ नियम $\equal[\eta]$ में कुछ नहीं जोड़ता। तब
  $\equal[\eta]$ में $\equal[\ext]$ समाहित है, और यदि $M \equal[\ext] N$,
  तो $M \equal[\eta] N$।

  यह सिद्ध करने के लिए कि $\equal[\eta]$, ~\ext{} के अधीन संवृत है, ध्यान
  दें कि \ext{} नियम द्वारा व्युत्पन्न किसी भी $M \equal N$ के लिए
  पूर्वाधार $Mx \equal[\eta] Nx$ है। तब $\equal$ के किसी नियम से
  $\lambd[x][Mx] \equal[\eta] \lambd[x][Nx]$; दोनों पक्षों पर $\eta$ लागू
  करने से $M \equal[\eta] N$ मिलता है।

  इसी प्रकार हम सिद्ध करते हैं कि $\eta$ नियम $\equal[ext]$ में समाहित है।
  $x \notin FV(M)$ वाले किन्हीं $\lambd[x][Mx]$ और $M$ के लिए
  $(\lambd[x][Mx])x \equal[\ext] Mx$; अतः \ext{} नियम से
  $\lambd[x][Mx] \equal[\ext] M$।
\end{proof}

\end{document}
